\documentclass[main.tex]{subfiles}
\begin{document}\newpage
\setdoublesep{0.35700 em}  % 'Bond Spacing'
\setatomsep{1.78500 em}    % 'Fixed Length'
\setbondoffset{0.18265 em} % 'Margin Width'
\newcommand{\bondwidth}{0.06642 em} % 'Line Width'
\setbondstyle{line width = \bondwidth}
\newgeometry{left=0.8in,right=0.8in, top=2.5cm,bottom=2cm}
\fancyhfoffset[E,O]{0pt}
\setlength{\columnsep}{30pt}
\begin{conclusion}
\end{conclusion}
\setstretch{0.3}




\begin{multicols*}{2}\setcounter{numA}{1}  %\newcounter{numA}


{\raggedright\textsc{\textbf{Electron-Dot structures of molecules }}\par}

%%%%%%%PROBLEM
\begin{question}[ID=\the\value{numA}]\SetQuestionProperties{section-title=\nameref{sec:units}}
Draw the electron-dot structure of:
\begin{inparaenum}[(a)]
\item  \ce{BH3}
\item  \ce{CH4}
\end{inparaenum}
\begin{center}\begin{tikzpicture}
\draw[step=.5cm,gray,very thin] (-4,-2) grid (2,2);
\end{tikzpicture}\end{center}
\end{question}
\begin{solution}\begin{inparaenum}[(a)]
\item  \chemfig{B( (-[:0]H)(-[:180]H)(-[:270]H))} 
\item  \chemfig{C( (-[:0]H)(-[:90]H)(-[:180]H)(-[:270]H))} 
\end{inparaenum}\hspace{0.1cm}\end{solution}
\stepcounter{numA}
%%%%%%%%%%%%%%









%%%%%%%PROBLEM
\begin{question}[ID=\the\value{numA}]\SetQuestionProperties{section-title=\nameref{sec:units}}
Draw the electron-dot structure of:
\begin{inparaenum}[(a)]
\item \ce{NH4^{+}}
\item \ce{H3O^{+}}
\end{inparaenum}
\begin{center}\begin{tikzpicture}
\draw[step=.5cm,gray,very thin] (-4,-2) grid (2,2);
\end{tikzpicture}\end{center}
\end{question}
\begin{solution} 
 \begin{inparaenum}[(a)]
 \item  \chemleft[\chemfig{N( (-[:0]H)(-[:90]H)(-[:180]H)(-[:270]H))}\chemright]\chemmove{\node[] at (0pt,35pt) {\footnotesize $+$};} 
 \item  \chemleft[\chemfig{N( (-[:0]H)(-[:90]H)(-[:180]H)(-[:270]H))}\chemright]\chemmove{\node[] at (0pt,35pt) {\footnotesize $+$};} 
 \end{inparaenum}\hspace{0.1cm}\end{solution}
\stepcounter{numA}
%%%%%%%%%%%%%%
%%%%%%%PROBLEM
\begin{question}[ID=\the\value{numA}]\SetQuestionProperties{section-title=\nameref{sec:units}}
The electron-dot structure of HI is:
\vspace{0.3cm}\begin{multicols*}{2}
\begin{enumerate}[label=(\alph*)]
\item \chemfig{ \lewis{46,H}-[:0]\lewis{026,I}}\hspace{.05in}
\item \chemfig{ \lewis{46,H}-[:0]\lewis{26,I}}\hspace{.05in}
\item \chemfig{ \lewis{46,H}-[:0]\lewis{6,I}}\hspace{.05in} 
\item \chemfig{ \lewis{26,H}-[:0]\lewis{26,I}}\hspace{.05in}
\item \chemfig{ H-[:0]\lewis{026,I}}\hspace{.05in}
\end{enumerate}\end{multicols*}
\end{question}
\begin{solution}
\chemfig{ H-[:0]\lewis{026,I}}\hspace{.05in}
\end{solution}
\stepcounter{numA}
%%%%%%%%%%%%%%
%%%%%%%PROBLEM
\begin{question}[ID=\the\value{numA}]\SetQuestionProperties{section-title=\nameref{sec:units}}
Indicate the charge of the central atom in the following electron-dot structure:
\vspace{0.3cm}\begin{multicols*}{2}
\begin{enumerate}[label=(\alph*)]
\item \chemleft[\chemfig{\lewis{2,C}( (-[:0]H)(-[:180]H)(-[:270]H))}\chemright] \chemmove{\node[] at (0pt,15pt) {\footnotesize $ $};}
\item \chemleft[\chemfig{\lewis{2,O}( (-[:0]H)(-[:180]H)(-[:270]H))}\chemright]
\chemmove{\node[] at (0pt,15pt) {\footnotesize $ $};}
\end{enumerate}\end{multicols*}
\end{question}
\begin{solution} 
 \begin{inparaenum}[(a)]
\item -2
\item +1
 \end{inparaenum}\hspace{0.1cm}\end{solution}
\stepcounter{numA}
%%%%%%%%%%%%%%
%%%%%%%PROBLEM
\begin{question}[ID=\the\value{numA}]\SetQuestionProperties{section-title=\nameref{sec:units}}
Indicate the charge of the central atom in the following electron-dot structure:
\vspace{0.3cm}\begin{multicols*}{2}
\begin{enumerate}[label=(\alph*)]
\item \chemleft[\chemfig{N( (-[:0]H)(-[:90]H)(-[:180]H)(-[:270]H))}\chemright]\chemmove{\node[] at (0pt,35pt) {\footnotesize $ $};}
\item \chemleft[\chemfig{C((-[:90]H)(-[:180]H)(-[:270]H))-[:0]\lewis{026,O}}\chemright]\chemmove{\node[] at (0pt,35pt) {\footnotesize $ $};}
\end{enumerate}
\end{multicols*}
\end{question}
\begin{solution} 
 \begin{inparaenum}[(a)]
\item +1
\item -1
 \end{inparaenum}\hspace{0.1cm}\end{solution}
\stepcounter{numA}
%%%%%%%%%%%%%%

{\raggedright\textsc{\textbf{Molecular shape}}\par}
%%%%%%%PROBLEM
\begin{question}[ID=\the\value{numA}]\SetQuestionProperties{section-title=\nameref{sec:units}}
Identify the molecular shape of the molecules:
\begin{inparaenum}[(a)]
\ce{NH3} %Trigonal pyramidal
\ce{CH4} %Tetrahedral
\end{inparaenum}
\end{question}
\begin{solution} 
 \begin{inparaenum}[(a)]
\ce{NH3}  (Trigonal pyramidal)
\ce{CH4}  (Tetrahedral)
 \end{inparaenum}\hspace{0.1cm}\end{solution}
\stepcounter{numA}
%%%%%%%%%%%%%%
%%%%%%%PROBLEM
\begin{question}[ID=\the\value{numA}]\SetQuestionProperties{section-title=\nameref{sec:units}}
Identify the molecular shape of the molecules:
\begin{inparaenum}[(a)]
\item \ce{H2}	%Linear
\item  \ce{BeCl2}	%Linear
  \ce{BF3} %Trigonal planar
\end{inparaenum}
\end{question}
\begin{solution} 
 \begin{inparaenum}[(a)]
\item \ce{H2}	 (Linear)
\item  \ce{BeCl2}	(Linear)
  \ce{BF3}  (Trigonal planar)
 \end{inparaenum}\hspace{0.1cm}\end{solution}
\stepcounter{numA}
%%%%%%%%%%%%%%



{\raggedright\textsc{\textbf{Polarity}}\par}

%%%%%%%PROBLEM
\begin{question}[ID=\the\value{numA}]\SetQuestionProperties{section-title=\nameref{sec:units}}
Indicate the polarity or non-polarity for the following molecules:
\begin{inparaenum}[(a)]
\item \ce{H2O}	%(Polar)
\item  \ce{HCl}	%(Polar)
\item  \ce{H2}	%(non Polar)
\end{inparaenum}
\end{question}
\begin{solution} 
 \begin{inparaenum}[(a)]
\item \ce{H2O}	%(Polar)
\item  \ce{HCl}	%(Polar)
\item  \ce{H2}	%(non Polar)
 \end{inparaenum}\hspace{0.1cm}\end{solution}
\stepcounter{numA}
%%%%%%%%%%%%%%

%%%%%%%PROBLEM
\begin{question}[ID=\the\value{numA}]\SetQuestionProperties{section-title=\nameref{sec:units}}
Indicate the polarity or non-polarity for the following molecules:
\begin{inparaenum}[(a)]
\item  \ce{NH3}	%(Polar)
\item  \ce{CO2}	%(non Polar)
\end{inparaenum}
\end{question}
\begin{solution} 
 \begin{inparaenum}[(a)]
\item  \ce{NH3}	%(Polar)
\item  \ce{CO2}	%(non Polar)
 \end{inparaenum}\hspace{0.1cm}\end{solution}
\stepcounter{numA}
%%%%%%%%%%%%%%



{\raggedright\textsc{\textbf{Hybrid Orbitals}}\par}
%%%%%%%PROBLEM
\begin{question}[ID=\the\value{numA}]\SetQuestionProperties{section-title=\nameref{sec:units}}
Indicate the hybridization of:
\begin{inparaenum}[(a)]
\item \ce{NH3}  %( $sp^3$)
\item \ce{CH4} %( $sp^3$)
\item \ce{H2O}%( $sp^3$)
\end{inparaenum}
\end{question}
\begin{solution} 
 \begin{inparaenum}[(a)]
\item \ce{NH3}   ( $sp^3$)
\item \ce{CH4}  ( $sp^3$)
\item \ce{H2O} ( $sp^3$)
 \end{inparaenum}\hspace{0.1cm}\end{solution}
\stepcounter{numA}
%%%%%%%%%%%%%%
%%%%%%%PROBLEM
\begin{question}[ID=\the\value{numA}]\SetQuestionProperties{section-title=\nameref{sec:units}}
Indicate the hybridization of:
\begin{inparaenum}[(a)]
\item \ce{NH3}  %( $sp^3$)
\item \ce{CH4} %( $sp^3$)
\item \ce{H2O}%( $sp^3$)
\end{inparaenum}
\end{question}
\begin{solution} 
 \begin{inparaenum}[(a)]
\item \ce{NH3}   ( $sp^3$)
\item \ce{CH4}  ( $sp^3$)
\item \ce{H2O} ( $sp^3$)
 \end{inparaenum}\hspace{0.1cm}\end{solution}
\stepcounter{numA}
%%%%%%%%%%%%%%

{\raggedright\textsc{\textbf{Molecular Orbital Theory}}\par}


%%%%%%%PROBLEM
\begin{question}[ID=\the\value{numA}]\SetQuestionProperties{section-title=\nameref{sec:units}}
Using the MO order provided below 
\[\sigma^{}_{2s}\sigma^{*}_{2s}\pi^{}_{2p}\sigma^{}_{2p}\pi^{*}_{2p}\sigma^{*}_{2p}\]
obtain the MO configuration for:
\begin{inparaenum}[(a)]
 \item \ce{B2} % ($\sigma^{2}_{2s}\sigma^{*2}_{2s}\pi^{2}_{2p}$)
 \item \ce{C2} % ($\sigma^{2}_{2s}\sigma^{*2}_{2s}\pi^{4}_{2p}$)
\end{inparaenum}
\end{question}
\begin{solution} 
 \begin{inparaenum}[(a)]
 \item \ce{B2}  ($\sigma^{2}_{2s}\sigma^{*2}_{2s}\pi^{2}_{2p}$)
 \item \ce{C2}  ($\sigma^{2}_{2s}\sigma^{*2}_{2s}\pi^{4}_{2p}$)
  \end{inparaenum}\hspace{0.1cm}\end{solution}
\stepcounter{numA}
%%%%%%%%%%%%%%
%%%%%%%PROBLEM
\begin{question}[ID=\the\value{numA}]\SetQuestionProperties{section-title=\nameref{sec:units}}
Using the MO order provided below 
\[\sigma^{}_{2s}\sigma^{*}_{2s}\pi^{}_{2p}\sigma^{}_{2p}\pi^{*}_{2p}\sigma^{*}_{2p}\]
obtain the MO configuration for:
\begin{inparaenum}[(a)]
 \item \ce{O2}  %($\sigma^{2}_{2s}\sigma^{2*}_{2s} \sigma^{2}_{2p}\pi^{4}_{2p}\pi^{2*}_{2p}$)
 \item \ce{F2^+}  %($\sigma^{2}_{2s}\sigma^{2*}_{2s} \sigma^{2}_{2p}\pi^{4}_{2p}\pi^{3*}_{2p}$)
\end{inparaenum}
\end{question}
\begin{solution} 
 \begin{inparaenum}[(a)]
 \item \ce{O2}  ($\sigma^{2}_{2s}\sigma^{2*}_{2s} \sigma^{2}_{2p}\pi^{4}_{2p}\pi^{2*}_{2p}$)
 \item \ce{F2^+}  ($\sigma^{2}_{2s}\sigma^{2*}_{2s} \sigma^{2}_{2p}\pi^{4}_{2p}\pi^{3*}_{2p}$)
 \end{inparaenum}\hspace{0.1cm}\end{solution}
\stepcounter{numA}
%%%%%%%%%%%%%%


%%%%%%%PROBLEM
\begin{question}[ID=\the\value{numA}]\SetQuestionProperties{section-title=\nameref{sec:units}}
Indicate the magnetic (paramagnetic or diamagnegtic) configuration of the molecule with MO configuration: $\sigma^{2}_{2s}\sigma^{2*}_{2s} \sigma^{2}_{2p}\pi^{4}_{2p}\pi^{3*}_{2p}$
\end{question}
\begin{solution} 
 paramagnetic\hspace{0.1cm}\end{solution}
\stepcounter{numA}
%%%%%%%%%%%%%%

 
\end{multicols*}
\newpage
\begin{answersenvironment}
\begin{minipage}[c]{1\textwidth}
\begin{localsize}{10}
{\Large \bf Answers}
\SetupExSheets{
  headings = inline-nr , % numbered and inline
  counter-format = qu) , % numbers 1) 2) ... 
}
%\printsolutions 
\printsolutions[byID={1,3,5,7,9,11,13}]
\end{localsize}
\end{minipage}\end{answersenvironment}

\end{document}















%
%
%
%%
%%\begin{center}\hspace{.05in}\chemfig{B(-[:0]\lewis{026,F})(-[:120]\lewis{024,F})(-[:240]\lewis{046,F})}\hspace{.05in}\hspace{.15in}\chemfig{\lewis{2,S}(=[:330]\lewis{06,O})(-[:210]\lewis{246,O})}\hspace{.05in} and \hspace{.05in}\chemfig{C(-[:0]H)(-[:90]H)(-[:180]H)(-[:270]H)}\hspace{.05in}
%%\end{center}
%
%\restoregeometry
%\end{enumerate}
%\end{multicols*}
%\pagecolor{green!10}\afterpage{\nopagecolor}\newpage
%\end{document}
